% =============================================================================
% ABSTRACT — Refined for clarity, accuracy, and impact
% =============================================================================
% Use % =============================================================================
% ABSTRACT — Refined for clarity, accuracy, and impact
% =============================================================================
% Use % =============================================================================
% ABSTRACT — Refined for clarity, accuracy, and impact
% =============================================================================
% Use % =============================================================================
% ABSTRACT — Refined for clarity, accuracy, and impact
% =============================================================================
% Use \input{paper_drafts/ABSTRACT} in main.tex, or copy the paragraph below.
%
% Plain-text version (no LaTeX):
%   Unsupervised physical parameter estimation from video lacks a common
%   benchmark: existing methods are evaluated on disjoint synthetic setups,
%   the only public real-world dataset is limited to single-body systems, and
%   there is no standard protocol for both equation identification and
%   parameter accuracy. We introduce IRIS, a benchmark of over 200 real-world
%   videos with high temporal resolution (60 fps), covering single-body
%   dynamics (e.g., pendulum, free fall, sliding, rotation) and multi-body
%   dynamics (e.g., hitting cones, two pendulums). Each setting is recorded
%   under controlled conditions with independently measured ground-truth
%   parameters and uncertainty estimates (mean, std, min, max), and each
%   phenomenon is paired with its governing ODE to support principled ablation.
%   We define a standardized evaluation protocol along five axes: parameter
%   accuracy, governing-equation selection accuracy, identifiability,
%   robustness, and extrapolation. We compare a latent-space physics baseline
%   (path-based equation selection, one-step loss) with a VLM-augmented
%   pipeline that identifies the equation family from video and estimates
%   parameters via a shared ODE bank and optional multi-step loss; we also
%   situate encoder-decoder and neural-implicit approaches relative to this
%   benchmark. The dataset, annotations (parameters.json), evaluation toolkit,
%   and baseline implementations will be publicly released.

\begin{abstract}
Unsupervised physical parameter estimation from video lacks a common benchmark: existing methods are evaluated on disjoint synthetic setups, the only public real-world dataset is limited to single-body systems, and there is no standard protocol for both equation identification and parameter accuracy. We introduce \textbf{IRIS}, a benchmark of over 200 real-world videos with high temporal resolution (60\,fps), covering single-body dynamics (e.g., pendulum, free fall, sliding, rotation) and multi-body dynamics (e.g., hitting cones, two pendulums). Each setting is recorded under controlled conditions with independently measured ground-truth parameters and uncertainty estimates (mean, std, min, max), and each phenomenon is paired with its governing ODE to support principled ablation. We define a standardized evaluation protocol along five axes: parameter accuracy, governing-equation selection accuracy, identifiability, robustness, and extrapolation. We compare a latent-space physics baseline (path-based equation selection, one-step loss) with a VLM-augmented pipeline that identifies the equation family from video and estimates parameters via a shared ODE bank and optional multi-step loss; we also situate encoder-decoder and neural-implicit approaches relative to this benchmark. The dataset, annotations (\texttt{parameters.json}), evaluation toolkit, and baseline implementations will be publicly released.
\end{abstract}
 in main.tex, or copy the paragraph below.
%
% Plain-text version (no LaTeX):
%   Unsupervised physical parameter estimation from video lacks a common
%   benchmark: existing methods are evaluated on disjoint synthetic setups,
%   the only public real-world dataset is limited to single-body systems, and
%   there is no standard protocol for both equation identification and
%   parameter accuracy. We introduce IRIS, a benchmark of over 200 real-world
%   videos with high temporal resolution (60 fps), covering single-body
%   dynamics (e.g., pendulum, free fall, sliding, rotation) and multi-body
%   dynamics (e.g., hitting cones, two pendulums). Each setting is recorded
%   under controlled conditions with independently measured ground-truth
%   parameters and uncertainty estimates (mean, std, min, max), and each
%   phenomenon is paired with its governing ODE to support principled ablation.
%   We define a standardized evaluation protocol along five axes: parameter
%   accuracy, governing-equation selection accuracy, identifiability,
%   robustness, and extrapolation. We compare a latent-space physics baseline
%   (path-based equation selection, one-step loss) with a VLM-augmented
%   pipeline that identifies the equation family from video and estimates
%   parameters via a shared ODE bank and optional multi-step loss; we also
%   situate encoder-decoder and neural-implicit approaches relative to this
%   benchmark. The dataset, annotations (parameters.json), evaluation toolkit,
%   and baseline implementations will be publicly released.

\begin{abstract}
Unsupervised physical parameter estimation from video lacks a common benchmark: existing methods are evaluated on disjoint synthetic setups, the only public real-world dataset is limited to single-body systems, and there is no standard protocol for both equation identification and parameter accuracy. We introduce \textbf{IRIS}, a benchmark of over 200 real-world videos with high temporal resolution (60\,fps), covering single-body dynamics (e.g., pendulum, free fall, sliding, rotation) and multi-body dynamics (e.g., hitting cones, two pendulums). Each setting is recorded under controlled conditions with independently measured ground-truth parameters and uncertainty estimates (mean, std, min, max), and each phenomenon is paired with its governing ODE to support principled ablation. We define a standardized evaluation protocol along five axes: parameter accuracy, governing-equation selection accuracy, identifiability, robustness, and extrapolation. We compare a latent-space physics baseline (path-based equation selection, one-step loss) with a VLM-augmented pipeline that identifies the equation family from video and estimates parameters via a shared ODE bank and optional multi-step loss; we also situate encoder-decoder and neural-implicit approaches relative to this benchmark. The dataset, annotations (\texttt{parameters.json}), evaluation toolkit, and baseline implementations will be publicly released.
\end{abstract}
 in main.tex, or copy the paragraph below.
%
% Plain-text version (no LaTeX):
%   Unsupervised physical parameter estimation from video lacks a common
%   benchmark: existing methods are evaluated on disjoint synthetic setups,
%   the only public real-world dataset is limited to single-body systems, and
%   there is no standard protocol for both equation identification and
%   parameter accuracy. We introduce IRIS, a benchmark of over 200 real-world
%   videos with high temporal resolution (60 fps), covering single-body
%   dynamics (e.g., pendulum, free fall, sliding, rotation) and multi-body
%   dynamics (e.g., hitting cones, two pendulums). Each setting is recorded
%   under controlled conditions with independently measured ground-truth
%   parameters and uncertainty estimates (mean, std, min, max), and each
%   phenomenon is paired with its governing ODE to support principled ablation.
%   We define a standardized evaluation protocol along five axes: parameter
%   accuracy, governing-equation selection accuracy, identifiability,
%   robustness, and extrapolation. We compare a latent-space physics baseline
%   (path-based equation selection, one-step loss) with a VLM-augmented
%   pipeline that identifies the equation family from video and estimates
%   parameters via a shared ODE bank and optional multi-step loss; we also
%   situate encoder-decoder and neural-implicit approaches relative to this
%   benchmark. The dataset, annotations (parameters.json), evaluation toolkit,
%   and baseline implementations will be publicly released.

\begin{abstract}
Unsupervised physical parameter estimation from video lacks a common benchmark: existing methods are evaluated on disjoint synthetic setups, the only public real-world dataset is limited to single-body systems, and there is no standard protocol for both equation identification and parameter accuracy. We introduce \textbf{IRIS}, a benchmark of over 200 real-world videos with high temporal resolution (60\,fps), covering single-body dynamics (e.g., pendulum, free fall, sliding, rotation) and multi-body dynamics (e.g., hitting cones, two pendulums). Each setting is recorded under controlled conditions with independently measured ground-truth parameters and uncertainty estimates (mean, std, min, max), and each phenomenon is paired with its governing ODE to support principled ablation. We define a standardized evaluation protocol along five axes: parameter accuracy, governing-equation selection accuracy, identifiability, robustness, and extrapolation. We compare a latent-space physics baseline (path-based equation selection, one-step loss) with a VLM-augmented pipeline that identifies the equation family from video and estimates parameters via a shared ODE bank and optional multi-step loss; we also situate encoder-decoder and neural-implicit approaches relative to this benchmark. The dataset, annotations (\texttt{parameters.json}), evaluation toolkit, and baseline implementations will be publicly released.
\end{abstract}
 in main.tex, or copy the paragraph below.
%
% Plain-text version (no LaTeX):
%   Unsupervised physical parameter estimation from video lacks a common
%   benchmark: existing methods are evaluated on disjoint synthetic setups,
%   the only public real-world dataset is limited to single-body systems, and
%   there is no standard protocol for both equation identification and
%   parameter accuracy. We introduce IRIS, a benchmark of over 200 real-world
%   videos with high temporal resolution (60 fps), covering single-body
%   dynamics (e.g., pendulum, free fall, sliding, rotation) and multi-body
%   dynamics (e.g., hitting cones, two pendulums). Each setting is recorded
%   under controlled conditions with independently measured ground-truth
%   parameters and uncertainty estimates (mean, std, min, max), and each
%   phenomenon is paired with its governing ODE to support principled ablation.
%   We define a standardized evaluation protocol along five axes: parameter
%   accuracy, governing-equation selection accuracy, identifiability,
%   robustness, and extrapolation. We compare a latent-space physics baseline
%   (path-based equation selection, one-step loss) with a VLM-augmented
%   pipeline that identifies the equation family from video and estimates
%   parameters via a shared ODE bank and optional multi-step loss; we also
%   situate encoder-decoder and neural-implicit approaches relative to this
%   benchmark. The dataset, annotations (parameters.json), evaluation toolkit,
%   and baseline implementations will be publicly released.

\begin{abstract}
Unsupervised physical parameter estimation from video lacks a common benchmark: existing methods are evaluated on disjoint synthetic setups, the only public real-world dataset is limited to single-body systems, and there is no standard protocol for both equation identification and parameter accuracy. We introduce \textbf{IRIS}, a benchmark of over 200 real-world videos with high temporal resolution (60\,fps), covering single-body dynamics (e.g., pendulum, free fall, sliding, rotation) and multi-body dynamics (e.g., hitting cones, two pendulums). Each setting is recorded under controlled conditions with independently measured ground-truth parameters and uncertainty estimates (mean, std, min, max), and each phenomenon is paired with its governing ODE to support principled ablation. We define a standardized evaluation protocol along five axes: parameter accuracy, governing-equation selection accuracy, identifiability, robustness, and extrapolation. We compare a latent-space physics baseline (path-based equation selection, one-step loss) with a VLM-augmented pipeline that identifies the equation family from video and estimates parameters via a shared ODE bank and optional multi-step loss; we also situate encoder-decoder and neural-implicit approaches relative to this benchmark. The dataset, annotations (\texttt{parameters.json}), evaluation toolkit, and baseline implementations will be publicly released.
\end{abstract}
